\hypertarget{openbmc-rest-api}{%
\section{OpenBMC REST API}\label{openbmc-rest-api}}

The OpenBMC REST API is disabled by default in bmcweb. Most of the
functionality is available on Redfish. This document provides some basic
structure and usage examples for the OpenBMC REST interface.

The schema for the rest interface is directly defined by the OpenBMC
D-Bus structure. Therefore, the objects, attributes and methods closely
map to those in the D-Bus schema.

For a quick explanation of HTTP verbs and how they relate to a RESTful
API, see \url{http://www.restapitutorial.com/lessons/httpmethods.html}.

\hypertarget{authentication}{%
\subsection{Authentication}\label{authentication}}

See the details on authentication at
\href{https://github.com/openbmc/docs/blob/master/REST-cheatsheet.md\#establish-rest-connection-session}{REST-cheatsheet}.

This tutorial uses token based authentication method:

\begin{verbatim}
export bmc=xx.xx.xx.xx
export token=`curl -k -H "Content-Type: application/json" -X POST https://${bmc}/login -d '{"username" : "root", "password" : "0penBmc"}' | grep token | awk '{print $2;}' | tr -d '"'`
curl -k -H "X-Auth-Token: $token" https://${bmc}/xyz/openbmc_project/...
\end{verbatim}

\hypertarget{http-get-operations-url-structure}{%
\subsection{HTTP GET operations \& URL
structure}\label{http-get-operations-url-structure}}

There are a few conventions on the URL structure of the OpenBMC rest
interface. They are:

\begin{itemize}
\tightlist
\item
  To query the attributes of an object, perform a GET request on the
  object name, with no trailing slash. For example:
\end{itemize}

\begin{verbatim}
curl -k -H "X-Auth-Token: $token" https://${bmc}/xyz/openbmc_project/inventory/system
\end{verbatim}

\begin{verbatim}
{
  "data": {
    "AssetTag": "",
    "BuildDate": "",
    "Cached": true,
    "FieldReplaceable": true,
    "Manufacturer": "",
    "Model": "",
    "PartNumber": "",
    "Present": true,
    "PrettyName": "",
    "SerialNumber": ""
  },
  "message": "200 OK",
  "status": "ok"
}
\end{verbatim}

\begin{itemize}
\tightlist
\item
  To query a single attribute, use the
  \texttt{attr/\textless{}name\textgreater{}} path. Using the
  \texttt{system} object from above, we can query just the \texttt{Name}
  value:
\end{itemize}

\begin{verbatim}
curl -k -H "X-Auth-Token: $token" https://${bmc}/xyz/openbmc_project/inventory/system/attr/Cached
\end{verbatim}

\begin{verbatim}
{
  "data": true,
  "message": "200 OK",
  "status": "ok"
}
\end{verbatim}

\begin{itemize}
\tightlist
\item
  When a path has a trailing-slash, the response will list the sub
  objects of the URL. For example, using the same object path as above,
  but adding a slash:
\end{itemize}

\begin{verbatim}
curl -k -H "X-Auth-Token: $token" https://${bmc}/xyz/openbmc_project/
\end{verbatim}

\begin{verbatim}
{
  "data": [
    "/xyz/openbmc_project/Chassis",
    "/xyz/openbmc_project/Hiomapd",
    "/xyz/openbmc_project/Ipmi",
    "/xyz/openbmc_project/certs",
    "/xyz/openbmc_project/console",
    "/xyz/openbmc_project/control",
    "/xyz/openbmc_project/dump",
    "/xyz/openbmc_project/events",
    "/xyz/openbmc_project/inventory",
    "/xyz/openbmc_project/ipmi",
    "/xyz/openbmc_project/led",
    "/xyz/openbmc_project/logging",
    "/xyz/openbmc_project/network",
    "/xyz/openbmc_project/object_mapper",
    "/xyz/openbmc_project/sensors",
    "/xyz/openbmc_project/software",
    "/xyz/openbmc_project/state",
    "/xyz/openbmc_project/time",
    "/xyz/openbmc_project/user"
  ],
  "message": "200 OK",
  "status": "ok"
}
\end{verbatim}

This shows that there are 19 children of the \texttt{openbmc\_project/}
object: \texttt{dump}, \texttt{software}, \texttt{control},
\texttt{network}, \texttt{logging},etc. This can be used with the base
REST URL (ie., \texttt{http://\$\{bmc\}/}), to discover all objects in
the hierarchy.

\begin{itemize}
\tightlist
\item
  Performing the same query with \texttt{/list} will list the child
  objects \emph{recursively}.
\end{itemize}

\begin{verbatim}
curl -k -H "X-Auth-Token: $token" https://${bmc}/xyz/openbmc_project/network/list
\end{verbatim}

\begin{verbatim}
{
  "data": [
    "/xyz/openbmc_project/network/config",
    "/xyz/openbmc_project/network/config/dhcp",
    "/xyz/openbmc_project/network/eth0",
    "/xyz/openbmc_project/network/eth0/ipv4",
    "/xyz/openbmc_project/network/eth0/ipv4/3b9faa36",
    "/xyz/openbmc_project/network/eth0/ipv6",
    "/xyz/openbmc_project/network/eth0/ipv6/ff81b6d6",
    "/xyz/openbmc_project/network/eth1",
    "/xyz/openbmc_project/network/eth1/ipv4",
    "/xyz/openbmc_project/network/eth1/ipv4/3b9faa36",
    "/xyz/openbmc_project/network/eth1/ipv4/66e63348",
    "/xyz/openbmc_project/network/eth1/ipv6",
    "/xyz/openbmc_project/network/eth1/ipv6/ff81b6d6",
    "/xyz/openbmc_project/network/host0",
    "/xyz/openbmc_project/network/host0/intf",
    "/xyz/openbmc_project/network/host0/intf/addr",
    "/xyz/openbmc_project/network/sit0",
    "/xyz/openbmc_project/network/snmp",
    "/xyz/openbmc_project/network/snmp/manager"
  ],
  "message": "200 OK",
  "status": "ok"
}
\end{verbatim}

\begin{itemize}
\tightlist
\item
  Adding \texttt{/enumerate} instead of \texttt{/list} will also include
  the attributes of the listed objects.
\end{itemize}

\begin{verbatim}
curl -k -H "X-Auth-Token: $token" https://${bmc}/xyz/openbmc_project/time/enumerate
\end{verbatim}

\begin{verbatim}
{
  "data": {
    "/xyz/openbmc_project/time/bmc": {
      "Elapsed": 1563209492098739
    },
    "/xyz/openbmc_project/time/host": {
      "Elapsed": 1563209492101678
    },
    "/xyz/openbmc_project/time/owner": {
      "TimeOwner": "xyz.openbmc_project.Time.Owner.Owners.BMC"
    },
    "/xyz/openbmc_project/time/sync_method": {
      "TimeSyncMethod": "xyz.openbmc_project.Time.Synchronization.Method.NTP"
    }
  },
  "message": "200 OK",
  "status": "ok"
}
\end{verbatim}

\hypertarget{http-put-operations}{%
\subsection{HTTP PUT operations}\label{http-put-operations}}

PUT operations are for updating an existing resource (an object or
property), or for creating a new resource when the client already knows
where to put it. These require a json formatted payload. To get an
example of what that looks like:

\begin{verbatim}
curl -k -H "X-Auth-Token: $token" https://${bmc}/xyz/openbmc_project/state/host0 > host.json
cat host.json
\end{verbatim}

\begin{verbatim}
{
  "data": {
    "AttemptsLeft": 3,
    "BootProgress": "xyz.openbmc_project.State.Boot.Progress.ProgressStages.Unspecified",
    "CurrentHostState": "xyz.openbmc_project.State.Host.HostState.Off",
    "OperatingSystemState": "xyz.openbmc_project.State.OperatingSystem.Status.OSStatus.Inactive",
    "RequestedHostTransition": "xyz.openbmc_project.State.Host.Transition.Off"
  },
  "message": "200 OK",
  "status": "ok"
}
\end{verbatim}

or

\begin{verbatim}
curl -k -H "X-Auth-Token: $token" https://${bmc}/xyz/openbmc_project/state/host0/attr/RequestedHostTransition > requested_host.json
cat requested_host.json
\end{verbatim}

\begin{verbatim}
{
  "data": "xyz.openbmc_project.State.Host.Transition.Off",
  "message": "200 OK",
  "status": "ok"
}
\end{verbatim}

When turning around and sending these as requests, delete the message
and status properties.

To make curl use the correct content type header use the -H option to
specify that we're sending JSON data:

\begin{verbatim}
curl -k -H "X-Auth-Token: $token" -H "Content-Type: application/json" -X PUT -d <json> <url>
\end{verbatim}

A PUT operation on an object requires a complete object. For partial
updates there is PATCH but that is not implemented yet. As a workaround
individual attributes are PUTable.

For example, make changes to the requested\_host.json file and do a PUT
(upload):

\begin{verbatim}
$ cat requested_host.json
{"data": "xyz.openbmc_project.State.Host.Transition.Off"}
$ curl -k -H "X-Auth-Token: $token" -H "Content-Type: application/json" -X PUT -T requested_host.json https://${bmc}/xyz/openbmc_project/state/host0/attr/RequestedHostTransition
\end{verbatim}

Alternatively specify the json inline with -d:

\begin{verbatim}
curl -k -H "X-Auth-Token: $token" -H "Content-Type: application/json" -X PUT -d '{"data": "xyz.openbmc_project.State.Host.Transition.On"}' https://${bmc}/xyz/openbmc_project/state/host0/attr/RequestedHostTransition
\end{verbatim}

When using `-d' just remember that json requires quoting.

\hypertarget{http-post-operations}{%
\subsection{HTTP POST operations}\label{http-post-operations}}

POST operations are for calling methods, but also for creating new
resources when the client doesn't know where to put it. OpenBMC does not
support creating new resources via REST so any attempt to create a new
resource will result in a HTTP 403 (Forbidden). These also require a
json formatted payload. To delete logging entries:

\begin{verbatim}
curl -k -H "X-Auth-Token: $token" -H 'Content-Type: application/json' -X POST -d '{"data":[]}' https://${bmc}/xyz/openbmc_project/logging/action/DeleteAll
\end{verbatim}

To invoke a method without parameters (Factory Reset of BMC and Host):

\begin{verbatim}
curl -k -H "X-Auth-Token: $token" -H 'Content-Type: application/json' -X POST -d '{"data":[]}' https://${bmc}/xyz/openbmc_project/software/action/Reset
\end{verbatim}

\hypertarget{http-delete-operations}{%
\subsection{HTTP DELETE operations}\label{http-delete-operations}}

DELETE operations are for removing instances. Only D-Bus objects
(instances) can be removed. If the underlying D-Bus object implements
the \texttt{xyz.openbmc\_project.Object.Delete} interface the REST
server will call it. If \texttt{xyz.openbmc\_project.Object.Delete} is
not implemented, the REST server will return a HTTP 403 (Forbidden)
error.

For example, to delete a event record:

Display logging entries:

\begin{verbatim}
curl -k -H "X-Auth-Token: $token" -H "Content-Type: application/json" -X GET https://${bmc}/xyz/openbmc_project/logging/entry/enumerate
\end{verbatim}

Then delete the event record with ID 1:

\begin{verbatim}
curl -k -H "X-Auth-Token: $token" -H "Content-Type: application/json" -X DELETE  https://${bmc}/xyz/openbmc_project/logging/entry/1
\end{verbatim}

\hypertarget{uploading-images}{%
\subsection{Uploading images}\label{uploading-images}}

It is possible to upload software upgrade images (for example to upgrade
the BMC or host software) via REST. The content-type should be set to
``application/octet-stream''.

For example, to upload an image:(the
\texttt{\textless{}file\_to\_upload\textgreater{}} must be a tar ball)

\begin{verbatim}
curl -k -H "X-Auth-Token: $token" -H "Content-Type: application/octet-stream" -X POST -T <file_to_upload> https://${bmc}/upload/image
\end{verbatim}

In above example, the filename on the BMC will be chosen by the REST
server.

It is possible for the user to choose the uploaded file's remote name:

\begin{verbatim}
curl -k -H "X-Auth-Token: $token" -H "Content-Type: application/octet-stream" -X PUT -T foo https://${bmc}/upload/image/bar
\end{verbatim}

In above example, the file foo will be saved with the name bar on the
BMC.

The operation will either return the version id (hash) of the uploaded
file on success:

\begin{verbatim}
{
  "data": "afb92384",
  "message": "200 OK",
  "status": "ok"
}
\end{verbatim}

or an error message

\begin{verbatim}
{
  "data": {
    "description": "Version already exists or failed to be extracted"
  },
  "message": "400 Bad Request",
  "status": "error"
}
\end{verbatim}

For more details on uploading and updating software, see:
\url{https://github.com/openbmc/docs/tree/master/architecture/code-update}

\hypertarget{event-subscription-protocol}{%
\subsection{Event subscription
protocol}\label{event-subscription-protocol}}

It is possible to subscribe to events, of interest, occurring on the
BMC. The implementation on the BMC uses WebSockets for this purpose, so
that clients don't have do employ polling. Instead, the rest server on
the BMC can push data to clients over a websocket. The BMC can push out
information pertaining to D-Bus InterfacesAdded and PropertiesChanged
signals.

Following is a description of the event subscription protocol, with
example JS code snippets denoting client-side code.

\begin{enumerate}
\def\labelenumi{\alph{enumi}.}
\item
  The client needs to have logged on to the BMC.
\item
  The client needs to open a secure websocket with the URL
  \texttt{\textless{}BMC\ IP\textgreater{}/subscribe}.
\end{enumerate}

\begin{verbatim}
var ws = new WebSocket("wss://<BMC IP>/subscribe");
\end{verbatim}

\begin{enumerate}
\def\labelenumi{\alph{enumi}.}
\setcounter{enumi}{2}
\tightlist
\item
  The client needs to send, over the websocket, a JSON dictionary,
  comprising of key-value pairs. This dictionary serves as the ``events
  filter''. All the keys are optional, so the dictionary can be empty if
  no filtering is desired. The filters represented by each of the
  key-value pairs are ORed.
\end{enumerate}

One of the supported keys is ``paths''. The corresponding value is an
array of D-Bus paths. The InterfacesAdded and PropertiesChanged D-Bus
signals emanating from any of these path(s) alone, and not from any
other paths, will be included in the event message going out of the BMC.

The other supported key is ``interfaces''. The corresponding value is an
array of D-Bus interfaces. The InterfacesAdded and PropertiesChanged
D-Bus signal messages comprising of any of these interfaces will be
included in the event message going out of the BMC.

All of the following are valid:

\begin{verbatim}
var data = JSON.stringify({
  paths: ["/xyz/openbmc_project/logging", "/xyz/openbmc_project/sensors"],
  interfaces: [
    "xyz.openbmc_project.Logging.Entry",
    "xyz.openbmc_project.Sensor.Value",
  ],
});
ws.onopen = function () {
  ws.send(data);
};
\end{verbatim}

\begin{verbatim}
var data = JSON.stringify({
  paths: ["/xyz/openbmc_project/logging", "/xyz/openbmc_project/sensors"],
});
ws.onopen = function () {
  ws.send(data);
};
\end{verbatim}

\begin{verbatim}
var data = JSON.stringify({
  interfaces: [
    "xyz.openbmc_project.Logging.Entry",
    "xyz.openbmc_project.Sensor.Value",
  ],
});
ws.onopen = function () {
  ws.send(data);
};
\end{verbatim}

\begin{verbatim}
var data = JSON.stringify({});
ws.onopen = function () {
  ws.send(data);
};
\end{verbatim}

\begin{enumerate}
\def\labelenumi{\alph{enumi}.}
\setcounter{enumi}{3}
\tightlist
\item
  The rest server on the BMC will respond over the websocket when a
  D-Bus event occurs, considering the client supplied filters. The rest
  servers notifies about InterfacesAdded and PropertiesChanged events.
  The response is a JSON dictionary as follows :
\end{enumerate}

InterfacesAdded

\begin{verbatim}
"event": InterfacesAdded
"path": <string : new D-Bus path that was created>
"interfaces": <dict : a dictionary of interfaces> (similar to org.freedesktop.DBus.ObjectManager.InterfacesAdded )
\end{verbatim}

PropertiesChanged

\begin{verbatim}
"event": PropertiesChanged
"path": <string : D-Bus path whose property changed>
"interface": <string : D-Bus interface to which the changed property belongs>
"properties": <dict : a dictionary of properties> (similar to org.freedesktop.DBus.Properties.PropertiesChanged)
\end{verbatim}
